% TOP_PROBLEM: {"id": 6600007, "title": "Coincidence Rank Conjecture", "category_id": 6, "category": {"id": 6, "name": "geometry", "display_name": "Geometry", "description": "Euclidean and non-Euclidean geometry, geometric structures.", "slug": "geometry", "order_index": 6, "created_at": "2026-07-31T15:26:25.671Z"}, "status": "partially_solved"}
% FIRST_POSTED: null
% ENTRY_KIND: statement_only
% Imported from: batch13.tex; UnsolvedMath ID 6600007
\documentclass[11pt]{article}
\usepackage[margin=1in]{geometry}
\usepackage{amsmath,amssymb,mathtools}
\usepackage{fontspec,unicode-math}
\setmainfont{Latin Modern Roman}
\setsansfont{Latin Modern Sans}
\setmathfont{Latin Modern Math}
\usepackage{xurl,hyperref}
\hypersetup{colorlinks=true,urlcolor=blue,linkcolor=blue}
\newcommand{\sref}[2]{\hyperlink{source:#1:#2}{[S#2]}}
\newcommand{\cataloguescope}[1]{\paragraph{Recorded mathematical question.}#1}
\begin{document}
% UnsolvedMath ID 6600007; source catalogue code AMR-065-0007
\setcounter{section}{6600006}
\section{Coincidence Rank Conjecture}\label{6600007}
{\small\sffamily UnsolvedMath ID 6600007 · Geometry}
\subsection{Definitions and mathematical statement}

\paragraph{Shared notation.}$\mathbb N=\{1,2,\ldots\}$, and $\mathbb Z,\mathbb Q,\mathbb R,\mathbb C$ denote the integers, rationals, reals and complex numbers. Any other conventions are specified locally.


A one-dimensional self-similar inflation tiling uses finitely many labeled interval types: expanding each interval by $\lambda$ and subdividing it according to a substitution reproduces the tiling. Its substitution (abelianisation) matrix $M$ counts tiles of each type in each substituted tile. The translation hull $\Omega$ is the closure of its translates in the local tiling topology, with the translation-invariant probability measure of the substitution system. Pure point dynamical spectrum means that $L^2(\Omega)$ is spanned by eigenfunctions of the translation action. A Pisot number is a real algebraic integer $\lambda>1$ whose other algebraic conjugates have modulus less than $1$. The maximal equicontinuous factor is the largest factor of the translation system with an equicontinuous action. The coincidence rank $c$ is the almost-everywhere multiplicity of the factor map $\Omega\to\Omega_{\rm eq}$. The algebraic norm $N(\lambda)$ is the product of all conjugates of $\lambda$, an integer; divisibility uses $|N(\lambda)|$.
\paragraph{Statement recorded in the supplied source.}
Does the coincidence rank of every one-dimensional Pisot inflation tiling divide $|N(\lambda)|$? \sref{6600007}{2}
\subsection{Short English statement}
Is the multiplicity of the maximal equicontinuous factor always a divisor of the inflation factor’s algebraic norm?
\subsection{Sources}
\begin{description}
\item[\hypertarget{source:6600007:1}{S1}] ULAM.AI and the UnsolvedMath Contributors, \emph{UnsolvedMath: A Curated Collection of Open Mathematics Problems}, record 6600007 (AMR-065-0007). CC BY 4.0. \url{https://huggingface.co/datasets/ulamai/UnsolvedMath}.
\item[\hypertarget{source:6600007:2}{S2}] Faustin Adiceam (compiler), \emph{Open Problems and Conjectures related to the Theory of Mathematical Quasicrystals}, arXiv:1604.06280v2 (2016), Section 2.2, Conjecture 2.2.3 and following definition, p. 4. \url{https://arxiv.org/pdf/1604.06280}.
\end{description}
\label{end:6600007}
\end{document}
